\section*{ข้อ 22 \hfill \normalsize MEC-0022}
(สอวน ม.5 2557) ระยะทาง $AB=D$ เป็นรางรถไฟตรง นกออกบินจาก $A$ พร้อมกันกับที่รถไฟออกจาก $B$ นกบินได้เร็ว $v$ รถไฟเคลื่อนที่เร็ว $u$ เมื่อนกถึงรถไฟครั้งแรกก็บินกลับทันทีด้วยความเร็วขนาดเท่าเดิม เมื่อถึงจุด $A$ รถไฟจะอยู่ห่าง $B$ เท่าใด (กำหนดว่า $v>u$ นกจะได้ไม่ถูกชน)
\begin{center}\begin{tikzpicture}[>=Stealth,scale=.8]\draw[dashed] (-.5,0)--(8,0);\draw[thick] (1,0)--(7,0);\fill (1,0) circle(2pt) node[below]{$A$};\fill (7,0) circle(2pt) node[below]{$B$};\node[above] at (1.7,.12){นก};\node[above] at (6.5,.12){รถไฟ};\draw[->,thick] (1.5,.55)--(3.1,.55) node[midway,above]{$v$};\draw[->,thick] (6.8,.55)--(5.2,.55) node[midway,above]{$u$};\end{tikzpicture}\end{center}
\medskip
